% !TEX root = ../main.tex
% Trazabilidad entre los cinco objetivos del enunciado original aprobado por la
% Comision de TFG y los seis objetivos especificos de esta memoria (IT-71).
%
% Por que hace falta: los criterios de evaluacion piden la correspondencia con
% «la propuesta del TFG original», que son esos cinco, y no con los OE1-OE6,
% que son una redaccion posterior. El mapeo era rastreable pero no estaba
% escrito en ninguna parte, de modo que reconstruirlo quedaba en manos del
% lector.
\begin{table}[htbp]
  \centering
  \small
  \caption[Trazabilidad con el enunciado original]{Correspondencia entre los
    objetivos del enunciado original del trabajo y los objetivos específicos de
    esta memoria. Fuente: Elaboración propia.}
  \label{tab:trazabilidad_enunciado}
  \begin{tabularx}{\textwidth}{@{}l >{\hsize=1.30\hsize}Y c >{\hsize=0.70\hsize}Y@{}}
    \toprule
    \textbf{N.\textsuperscript{o}} & \textbf{Objetivo del enunciado original} & \textbf{OE} & \textbf{Dónde se responde} \\
    \midrule
    1 & Desarrollar un rastreador de la web del centro que recoja titulaciones, guías docentes y memorias de verificación & OE1 & Apartado~\ref{secc:resultados_coleccion} \\
    2 & Estudiar los sistemas de generación aumentada por recuperación                                                    & OE2 · OE3 & Apartados~\ref{secc:resultados_fragmentacion} y~\ref{secc:resultados_embeddings} \\
    3 & Configurar y desplegar el sistema RAG sobre la colección construida                                               & OE4 & Apartado~\ref{secc:resultados_rag_sistema} \\
    4 & Construir una aplicación web de conversación sobre ese sistema                                                     & OE6 & Apartado~\ref{secc:fase3_interfaz} \\
    5 & Probar y validar el resultado                                                                                      & OE5 & Apartado~\ref{secc:resultados_rag_sistema} \\
    \bottomrule
  \end{tabularx}
  \begin{minipage}{0.95\textwidth}
    \vspace{4pt}
    \tiny
    \textbf{Notas:} \textbf{OE} remite al objetivo específico correspondiente
    del Capítulo~\ref{cap:objetivos}. La correspondencia no es uno a uno en los
    dos sentidos: el segundo objetivo del enunciado se desdobla en dos objetivos
    específicos porque la fragmentación y la elección del modelo de
    incrustaciones se decidieron con experimentos distintos.
  \end{minipage}
\end{table}
